\chapter{Fundamentos e revisão da literatura}
\label{cap:fundamentos}

Este capítulo reúne o que é necessário para ler os capítulos seguintes:
as polarizações na RG e em teorias métricas gerais, a gravitação
massiva linear e seus problemas conhecidos, os limites observacionais
sobre a massa do gráviton e o princípio de detecção por PTAs. A última
seção situa a dissertação na literatura recente.

\section{Ondas gravitacionais na relatividade geral linearizada}

Adota-se a assinatura $\eta_{\mu\nu}=\operatorname{diag}(1,-1,-1,-1)$,
como no TG. Longe das fontes, escreve-se
$g_{\mu\nu}=\eta_{\mu\nu}+h_{\mu\nu}$, com $|h_{\mu\nu}|\ll1$, e
define-se a perturbação com traço invertido
$\bar h_{\mu\nu}=h_{\mu\nu}-\frac12\eta_{\mu\nu}h$, onde
$h=\eta^{\mu\nu}h_{\mu\nu}$. Na condição harmônica
$\partial^\mu\bar h_{\mu\nu}=0$, as equações de Einstein no vácuo
reduzem-se a $\Box\bar h_{\mu\nu}=0$ \autocite{maggiore2008}. Uma onda
plana $\bar h_{\mu\nu}\propto e^{-i\omega t+ikz}$ tem $\omega=ck$.

A condição harmônica não fixa completamente o calibre. As
transformações residuais permitem anular as componentes temporais e o
traço, e impor transversalidade; no calibre transverso e sem traço (TT)
restam apenas $h_{11}=-h_{22}=h_+$ e $h_{12}=h_\times$. Esse é o
conteúdo físico da RG: duas polarizações, de helicidades $\pm2$
\autocite{maggiore2008}. O TG reproduz essa contagem no seu Capítulo~3.

Um detector não mede $h_{\mu\nu}$ diretamente, pois $h_{\mu\nu}$
depende do calibre. O observável local é a curvatura. Para duas
partículas livres próximas, separadas por $\xi^j$, a equação do desvio
geodésico dá a aceleração relativa em termos das componentes
\begin{equation}
 E_{ij}\equiv R_{0i0j},
 \label{eq:fund-marés}
\end{equation}
chamadas componentes ``elétricas'' do tensor de Riemann ou matriz de
marés. A convenção de sinal adotada e a relação
$\ddot\xi^i=c^2E_{ij}\xi^j$ estão no Apêndice~\ref{ap:convencoes}. Como
o Riemann linearizado é invariante por transformações infinitesimais de
coordenadas, $E_{ij}$ não depende do calibre. É ele que o TG calcula e é
dele que partirá a resposta de PTA no Capítulo~\ref{cap:pta}.

\section{Polarizações em teorias métricas}

A matriz $E_{ij}$ é simétrica e tem seis componentes. Para uma onda que
se propaga ao longo de $z$, elas se organizam por seu comportamento sob
rotações em torno do eixo de propagação:
\begin{itemize}
 \item dois modos tensoriais, $E_{11}-E_{22}$ e $E_{12}$, de helicidade
 $\pm2$, os únicos presentes na RG;
 \item dois modos vetoriais, $E_{13}$ e $E_{23}$, de helicidade $\pm1$;
 \item dois modos escalares, de helicidade zero: o transversal ou
 \emph{breathing}, $E_{11}+E_{22}$, e o longitudinal, $E_{33}$.
\end{itemize}
Uma teoria métrica genérica pode excitar até esses seis padrões. A
teoria determina quantos deles variam de modo independente e com que
relação de dispersão.

Eardley, Lee e Lightman organizaram essa descrição no formalismo de
Newman--Penrose \autocite{eardley1973}. Para ondas planas nulas, com
vetor de onda $k^\mu$ tal que $k^\mu k_\mu=0$, as identidades de Bianchi
reduzem as componentes independentes do Riemann a seis, descritas pelas
amplitudes $\Psi_2$ (real), $\Psi_3$ (complexa), $\Psi_4$ (complexa) e
$\Phi_{22}$ (real), com helicidades $0$, $\pm1$, $\pm2$ e $0$. O
subgrupo de Lorentz que preserva um vetor nulo é isomorfo ao grupo
euclidiano E(2). Sob suas transformações, a presença de algumas
amplitudes depende do observador, mas certas afirmações são invariantes.
Isso define as classes E(2):
\begin{itemize}
 \item II$_6$: $\Psi_2\neq0$; todos os observadores concordam sobre
 $\Psi_2$, e até seis modos podem aparecer;
 \item III$_5$: $\Psi_2=0$ e $\Psi_3\neq0$;
 \item N$_3$: $\Psi_2=\Psi_3=0$, com $\Psi_4\neq0$ e $\Phi_{22}\neq0$;
 \item N$_2$: apenas $\Psi_4\neq0$, como na RG;
 \item O$_1$: apenas $\Phi_{22}\neq0$; O$_0$: nenhuma onda.
\end{itemize}
Teorias escalares--tensoriais do tipo Brans--Dicke são da classe
N$_3$; a RG é N$_2$ \autocite{will2014}.

O ponto relevante para esta dissertação é o domínio da classificação.
Ela foi construída para ondas exatamente nulas e estendida a ondas
\emph{quase} nulas. Nesse caso, introduz-se
\begin{equation}
 \epsilon=\left(\frac{c}{v_g}\right)^2-1
 \label{eq:fund-epsilon}
\end{equation}
e desprezam-se correções de ordem $\epsilon$ nas relações entre as
componentes \autocite{eardley1973}. O TG e o artigo de 2004 adotam essa
extensão (TG, eqs.~5.7 e 5.33--5.43). Para uma onda massiva, $\epsilon$
não é necessariamente pequeno. O Capítulo~\ref{cap:polarizacoes} examina
esse ponto quantitativamente.

\section{Gravitação massiva linear}
\label{sec:fund-massiva}

\subsection{Uma família de termos de massa}

Em torno de Minkowski, a ação quadrática de um campo simétrico
$h_{\mu\nu}$ que reproduz a RG linearizada sem massa admite um termo de
massa com duas estruturas possíveis, $h_{\mu\nu}h^{\mu\nu}$ e $h^2$.
Nas convenções do Capítulo~\ref{cap:polarizacoes}, as equações no vácuo
podem ser escritas como
\begin{equation}
 G^{(1)}_{\mu\nu}[h]-\frac{\mu^2}{2}
 \left(h_{\mu\nu}-a\,\eta_{\mu\nu}h\right)=0,
 \qquad \mu=\frac{m_gc}{\hbar},
 \label{eq:fund-familia}
\end{equation}
em que $G^{(1)}_{\mu\nu}$ é o tensor de Einstein linearizado e $a$ é um
número real. A identidade de Bianchi linearizada,
$\partial^\mu G^{(1)}_{\mu\nu}\equiv0$, impõe quatro vínculos,
$\partial^\mu(h_{\mu\nu}-a\eta_{\mu\nu}h)=0$, que não são escolhas de
calibre: a massa quebra a invariância de calibre da teoria sem massa.
Dois casos importam aqui.
\begin{itemize}
 \item $a=1$ é o termo de Fierz--Pauli \autocite{fierz1939}. O traço da
 equação, combinado com os vínculos, dá $h=0$; restam
 $\partial^\mu h_{\mu\nu}=0$ e $(\Box+\mu^2)h_{\mu\nu}=0$. São
 $10-4-1=5$ graus de liberdade, as cinco helicidades de uma partícula
 massiva de spin dois.
 \item $a=\frac12$ é o termo de Visser. Os vínculos tornam-se
 $\partial^\mu\bar h_{\mu\nu}=0$, e a equação reduz-se a
 $(\Box+\mu^2)\bar h_{\mu\nu}=0$. O traço não é eliminado, e restam
 seis graus de liberdade.
\end{itemize}
O segundo caso é o do TG. O termo de massa de Visser é
\autocite{visser1998}
\begin{equation}
 \mathcal L_{\rm mass}\propto m_g^2\left[
 (g-g_0)_{\mu\nu}(g-g_0)^{\mu\nu}-\frac12(g-g_0)^2\right],
 \label{eq:fund-visser}
\end{equation}
com índices elevados pela métrica de fundo $g_0$. Linearizado com
$g_0=\eta$, isso corresponde a $a=\frac12$. Visser nota explicitamente
que esse termo não é o de Fierz--Pauli. Afirma também que a escolha é
essencial para um limite clássico bem comportado quando $m_g\to0$
\autocite{visser1998}.

\subsection{O sexto modo é um fantasma}
\label{sec:fund-fantasma}

Para $a\neq1$, o modo adicional é um escalar associado ao traço de
$h_{\mu\nu}$. Sua energia cinética tem sinal oposto ao dos demais
modos, isto é, trata-se de um fantasma. Essa é uma propriedade conhecida
dos termos de massa lineares que não são de Fierz--Pauli
\autocite{hinterbichler2012}. Park estudou explicitamente essa família.
Ele decompôs o campo em uma parte de spin dois, de massa $m$, e um
escalar fantasma de massa $m_a^2=(4a-1)m^2/[2(1-a)]$ \autocite{park2010}.
Para $a=\frac12$, $m_a=m$: o fantasma tem a mesma massa e, portanto, a
mesma relação de dispersão dos demais modos. Isso é consistente com a
equação única $(\Box+\mu^2)\bar h_{\mu\nu}=0$ do modelo de Visser.

O mesmo fantasma explica por que o modelo de Visser não sofre a
descontinuidade vDVZ na ordem linear. Com fontes conservadas, a solução
de $(\Box+\mu^2)\bar h_{\mu\nu}=-16\pi GT_{\mu\nu}/c^4$ tem a mesma
estrutura tensorial da RG, $T'^{\mu\nu}T_{\mu\nu}-\frac12T'T$,
multiplicada por um propagador de Yukawa. Em Fierz--Pauli, a estrutura
é $T'^{\mu\nu}T_{\mu\nu}-\frac13T'T$, e a diferença não desaparece
quando $m_g\to0$ \autocite{vandam1970,zakharov1970}. Em Fierz--Pauli, a
helicidade zero acopla-se ao traço de $T_{\mu\nu}$; no modelo de Visser,
o fantasma escalar acopla-se ao traço com sinal oposto e cancela esse
efeito a distâncias muito menores que $\mu^{-1}$ \autocite{park2010}.
O ``sexto estado de polarização'' do TG é, portanto, o grau de
liberdade que torna o modelo livre de vDVZ na ordem linear, ao preço de
uma instabilidade.

\subsection{Estrutura não linear}
\label{sec:fund-naolinear}

A consistência de Fierz--Pauli é linear. Boulware e Deser mostraram que
extensões não lineares genéricas reintroduzem um sexto modo fantasma
\autocite{boulware1972}. A construção de de Rham, Gabadadze e Tolley
(dRGT) escolhe o potencial de modo a preservar o vínculo que o elimina
\autocite{derham2011}; a versão bimétrica de Hassan e Rosen torna
dinâmica também a segunda métrica \autocite{hassan2012}. Em torno de
fundos planos, o setor massivo dessas teorias tem o conteúdo linear de
Fierz--Pauli; na versão bimétrica, ele vem acompanhado de um gráviton
sem massa \autocite{derham2014review}.

A descontinuidade vDVZ é resolvida por outro mecanismo. Perto de uma
fonte de raio de Schwarzschild $r_s$, as interações não lineares da
helicidade zero dominam abaixo do raio de Vainshtein,
$r_V\sim(r_s/\mu^2)^{1/3}$, e suprimem a força escalar adicional
\autocite{vainshtein1972,babichev2013}. Essa supressão tem consequência
direta para esta dissertação. A amplitude com que a helicidade zero é
\emph{emitida} por fontes astrofísicas depende da dinâmica não linear.
Não é fixada pela análise de propagação linear. Por isso, nos
Capítulos~\ref{cap:pta} e~\ref{cap:simulacoes}, a potência do modo
escalar será um parâmetro livre.

\section{Limites observacionais sobre a massa do gráviton}

A Tabela~\ref{tab:fund-limites} resume limites representativos, com a
frequência de corte $f_g=m_gc^2/h$ correspondente. Os métodos diferem.
Os limites por ondas gravitacionais testam a dispersão ao longo da
propagação. Os de efemérides planetárias testam um potencial de Yukawa
estático, cuja relação com a massa de dispersão depende da teoria e do
mecanismo de Vainshtein. Os valores não são, portanto, estritamente
intercambiáveis \autocite{will2014,lvk2026}.

\begin{table}[htbp]
\centering
\caption{Limites superiores representativos sobre a massa do gráviton e
frequência de corte correspondente, $f_g=m_gc^2/h$.}
\label{tab:fund-limites}
\small
\begin{tabular}{p{0.46\textwidth}cc}
\toprule
Origem & $m_g$ (eV$/c^2$) & $f_g$ (Hz)\\
\midrule
Sistema Solar, adotado no artigo de 2004 \autocite{depaula2004} & $4{,}4\times10^{-22}$ & $1{,}1\times10^{-7}$\\
GWTC-3, priori uniforme na massa, 90\% \autocite{lvk2026} & $2{,}23\times10^{-23}$ & $5{,}4\times10^{-9}$\\
GWTC-4, 90\% \autocite{lvk2026} & $1{,}92\times10^{-23}$ & $4{,}6\times10^{-9}$\\
NANOGrav 15 anos com modelo espectral \autocite{wu2024} & $8{,}2\times10^{-24}$ & $2{,}0\times10^{-9}$\\
Efemérides planetárias, 99,7\% (citado em \autocite{lvk2026}) & $1{,}01\times10^{-24}$ & $2{,}4\times10^{-10}$\\
\bottomrule
\end{tabular}
\end{table}

Duas observações orientam a leitura. Primeiro, todos os limites por
ondas gravitacionais situam $f_g$ em poucos nanohertz ou abaixo, isto é,
na banda das PTAs ou abaixo dela. Segundo, a análise de Wu e
colaboradores mostra que as correlações angulares do NANOGrav, sozinhas,
não fornecem um limite superior explícito no intervalo
$3\times10^{-25}$--$8\times10^{-24}\,\mathrm{eV}/c^2$. O limite da
tabela depende da combinação com o espectro observado \autocite{wu2024}.
Esse resultado antecipa um tema do Capítulo~\ref{cap:simulacoes}: a
informação sobre a massa contida na forma das correlações é pequena.

\section{Redes de temporização de pulsares}

Pulsares de milissegundo são relógios de grande estabilidade. Uma PTA
mede, para cada pulsar, os tempos de chegada (TOAs) dos pulsos ao longo
de anos, e subtrai um modelo de temporização que inclui rotação,
astrometria e órbita. Uma onda gravitacional que atravessa a linha de
visada altera a frequência aparente dos pulsos. O resíduo de
temporização é a integral temporal desse desvio fracionário de
frequência, $r_a(t)=\int^t z_a(t')\,dt'$ \autocite{detweiler1979}.

Para um fundo estocástico, isotrópico e não polarizado da RG, a
correlação entre os resíduos de dois pulsares separados por um ângulo
$\zeta$ tem a forma de Hellings--Downs \autocite{hellings1983},
\begin{equation}
 \Gamma_{\rm HD}(\zeta)=\frac12-\frac{x}{4}+\frac32x\ln x,
 \qquad x=\frac{1-\cos\zeta}{2}.
 \label{eq:fund-hd}
\end{equation}
Essa curva é a assinatura de uma onda gravitacional tensorial e sem
massa. Sua observação, com significância crescente, é a base da
evidência reportada em 2023 \autocite{agazie2023}. O espectro de um fundo
de binárias de buracos negros supermassivos é uma lei de potência,
$h_c(f)\propto f^{-2/3}$, que nos resíduos corresponde a um índice
espectral $\gamma=13/3$. Com observações de duração $T$, a resolução em
frequência é $1/T$. Para $T$ entre cinco e vinte anos, os canais mais
baixos ficam entre cerca de $1,6$ e $6\,\mathrm{nHz}$, a mesma ordem das
frequências de corte da Tabela~\ref{tab:fund-limites}.

\section{PTAs e gravidade massiva: estado da arte}

Lee, Jenet e Price calcularam as correlações de PTA para polarizações
não einsteinianas \autocite{lee2008}. Lee e colaboradores incluíram a
dispersão de um gráviton massivo na correlação tensorial, com o termo da
Terra, e estimaram a sensibilidade de PTAs à massa \autocite{lee2010}.
Liang e Trodden derivaram as correlações de todas as polarizações de
uma teoria massiva, com termos da Terra e do pulsar
\autocite{liang2021}. A revisão mais recente desse trabalho inclui o
movimento dos observadores quando a perturbação tem componentes
temporais. Cordes e colaboradores trataram distâncias finitas e
velocidades de fase arbitrárias, com expressões fechadas para o termo
da Terra e para a autocorrelação \autocite{cordes2025}. Bernardo e Ng
estudaram correlações escalares e tensoriais em teorias modificadas e
seu limite de variância cósmica \autocite{bernardo2024}.

No lado observacional, Wu e colaboradores limitaram a massa com os dados
de 15 anos do NANOGrav \autocite{wu2024}. A colaboração NANOGrav buscou
polarizações transversais adicionais sem encontrar evidência
significativa \autocite{nanograv2024}. Choi e Kahniashvili compararam
setores massivos com correlações binadas de diferentes PTAs
\autocite{choi2025}; Wu, Bi e Huang estudaram o efeito da interferência
de um número finito de fontes \autocite{wu2025}. Zhao e Wang analisaram
correlações comprimidas em frequência do MeerKAT. Eles avaliaram a
resposta dispersiva em uma frequência de referência $f=1/T$ e
projetaram limites para o SKA \autocite{zhao2026}. Han e Zhao
estudaram previsões com a covariância completa de estimadores de
correlação \autocite{han2026}. Franciolini e colaboradores discutiram as
verossimilhanças apropriadas para dados de fundos estocásticos
\autocite{franciolini2025}.

Do lado da cinemática local, Hyun, Kim e Lee obtiveram as amplitudes
exatas dos seis modos para propagação não nula. Eles apontaram as
limitações das aproximações quase nulas \autocite{hyun2019}.

\section{Onde esta dissertação se situa}

A literatura de PTAs trata as polarizações adicionais, em geral, como
setores fenomenológicos independentes, cada um com sua potência. A
literatura sobre o TG e seus sucessores trata a cinemática local, sem
chegar ao observável de temporização. Esta dissertação faz a ligação
entre as duas. O Capítulo~\ref{cap:polarizacoes} revisa a cinemática do
TG de forma exata e a reinterpreta. O Capítulo~\ref{cap:pta} mostra que
a resposta de PTA depende apenas da matriz de marés calculada no TG, e
calcula as correlações do setor de helicidade zero de Fierz--Pauli, no
qual as deformações transversal e longitudinal vêm de uma única
amplitude. O Capítulo~\ref{cap:simulacoes} avalia quanto disso uma rede
pequena consegue medir.
